\documentclass[11pt]{article}
\usepackage[a4paper,margin=26mm]{geometry}
\usepackage{amsmath,amssymb,amsthm,booktabs,array}
\usepackage[T1]{fontenc}
\usepackage[colorlinks=true,allcolors=blue,breaklinks=true]{hyperref}
\newtheorem{theorem}{Theorem}
\newtheorem{lemma}[theorem]{Lemma}
\newtheorem{corollary}[theorem]{Corollary}
\newcommand{\ket}[1]{\lvert #1\rangle}
\newcommand{\bra}[1]{\langle #1\rvert}
\newcommand{\tr}{\operatorname{tr}}
\title{Flat averages refute Zhan--Gupta--Hillery,\
Section 5, Conjecture 2 (2024, p.~10:11):\
an extension of Skotiniotis et al.\ (2018)}
\author{Haobo Ma\\\small ChronoAI Pte Ltd, Singapore
\and Wenlin Zhang\\\small National University of Singapore\\\small ChronoAI Pte Ltd, Singapore}
\date{November 5, 2026}
\begin{document}
\maketitle
\begin{abstract}
For Zhan--Gupta--Hillery, Section 5, Conjecture 2 (2024, p.~10:11;
arXiv:2306.17401v6, p.~11), we derive a strict-averaging counterexample
from the complementary-layer Gram formula of Skotiniotis et al.\ (2018).
Two different coefficient-square distributions and their whole averaging
segment have the same Gram matrix and, for three sensors at
$\theta=\pi/4$, the same error $(14-4\sqrt{10})/27$.
We extend the earlier symmetric-probe calculation by proving weak convexity
and the existence of a symmetric pure optimum without a symmetry assumption.
This proves the companion Conjecture 1 and answers the source's final question:
one probe is optimal for both minimum-error and unambiguous discrimination
at every detector angle. The reductions also apply to fixed unequal priors
and diagonal impact patterns, without solving their general optimizers.
\end{abstract}
\noindent\small Licensed under Creative Commons Attribution 4.0 International
(CC BY 4.0).\normalsize

\section{The posed statement and its Gram matrix}
The Gram formula and symmetric-probe optimization used below already appear
in Skotiniotis, Hotz, Calsamiglia and Mu\~noz-Tapia
(2018, Eq.~(8))~\cite{SHCM18}; their expanded article with Llorens
(2024, Section III.B, Eqs.~(12)--(17))~\cite{SLHCM24}
explicitly leaves the restriction to pure permutation-invariant probes
unproved. We use that formula to refute the later strict-averaging
conjecture and supply a direct argument removing the restriction.
Their phase parameter is $\varphi=2\theta$ (folded to $[0,\pi]$ if needed);
their 2024 Gram matrix includes the prior factor $1/n$, whereas ours
has diagonal one. The complementary-weight symmetry gives the flat family;
weak convexity supplies the pure-probe reduction. The earlier detector-sensor
model is described in~\cite{HGZ}.

Zhan, Gupta and Hillery~\cite{ZGH} consider $n\geq3$ qubit sensors,
one of which is affected with equal prior probability. In the tensor
eigenbasis of $U$, write
$U\ket{u_\pm}=e^{\pm i\theta}\ket{u_\pm}$, where $0<\theta<\pi$.
We use radians. A normalized pure probe $\ket\psi$ produces states
$\ket{\phi_i}=U_i\ket\psi$, $1\leq i\leq n$, where $U_i$ applies
$U$ to sensor $i$. Its minimum error is
\begin{equation}\label{eq:error}
 P(\psi)=1-\max_{M_i\succeq0,\,\sum_iM_i=I}
 \frac1n\sum_{i=1}^n\bra{\phi_i}M_i\ket{\phi_i}.
\end{equation}
The maximum exists: the POVM set in this finite-dimensional space is compact.
Conjecture 2 is at Section 5, p.~10:11 (arXiv v6, p.~11;
source label \texttt{conj:avg}). Its strict conclusion is that the
average ``has a lower objective function value.'' Thus different coefficient-squares and
$P(\psi)=P(\psi')$ imply $P(\psi_{\rm avg})<P(\psi')$, where
\begin{equation}\label{eq:average}
 \ket{\psi_{\rm avg}}=\sum_{x\in\{0,1\}^n}
 \sqrt{\frac{|\psi_x|^2+|\psi'_x|^2}{2}}\ket x.
\end{equation}
This is a new pure state, not a density-matrix mixture.

Put $p_x=|\psi_x|^2$ and $s_i(x)=2x_i-1$. The final-state Gram matrix is
\begin{equation}\label{eq:gram}
 G(p)_{ij}=\langle\phi_i\mid\phi_j\rangle
 =\sum_x p_x e^{i\theta(s_j(x)-s_i(x))}.
\end{equation}
Thus phases of the probe coefficients do not affect discrimination, and
$G(p)$ is affine in $p$.

\begin{lemma}\label{lem:gram}
Two labelled pure-state ensembles with identical Gram matrices have
the same attainable measurement probabilities, including the same optimal
minimum-error success for any fixed priors.
\end{lemma}
\begin{proof}
Send each labelled vector of the first ensemble to its counterpart in the
second. For any coefficient vector $z$, the squared norm of its linear
combination is $z^*Gz$ in both ensembles. The map is therefore well-defined
and isometric on their spans. Compress a POVM to a span and transport it by
this isometry. Extend it to the full ambient space by distributing the
orthogonal-complement projector among the outcomes. This preserves every
probability on the ensemble. Apply the inverse isometry for the converse.
\end{proof}

\section{An explicit flat family}
For $0\leq k\leq n$, let
\[
 \ket{D_k}=\binom nk^{-1/2}\sum_{|x|=k}\ket x,
 \qquad p^{(k)}_x=\binom nk^{-1}{\bf1}_{|x|=k}.
\]
These are Dicke states in the eigenbasis of $U$, as in~\cite{ZGH}.

\begin{theorem}\label{thm:flat}
Let $n\geq3$, $1\leq k\leq n-1$, $k\ne n-k$, and $0<\theta<\pi$.
For $0\leq t\leq1$ define
$\ket{A_t}=\sum_x\sqrt{t p^{(k)}_x+(1-t)p^{(n-k)}_x}\ket x$.
The coefficient-square distributions of $D_k$ and $D_{n-k}$ are distinct,
but
\[
 P(D_k)=P(D_{n-k})=P(A_t)\quad(0\leq t\leq1).
\]
In particular, $t=1/2$ refutes Conjecture 2.
\end{theorem}
\begin{proof}
For distinct sensors $i,j$, the number of weight-$k$ words with
$(x_i,x_j)=(1,0)$ equals the number with $(0,1)$, namely
$\binom{n-2}{k-1}$. The other
$\binom{n-2}{k}+\binom{n-2}{k-2}$ words have equal selected bits.
Equation~\eqref{eq:gram} consequently has diagonal $1$ and common
off-diagonal entry
\begin{equation}\label{eq:ck}
 c_k=1-\frac{4k(n-k)}{n(n-1)}\sin^2\theta.
\end{equation}
This is~\cite[Eq.~(8)]{SHCM18}, equivalently~\cite[Eq.~(14)]{SLHCM24}.
Here $\binom{n-2}{k-1}/\binom nk=k(n-k)/(n(n-1))$;
out-of-range binomial coefficients are zero.
Since $c_k=c_{n-k}$, the two Gram matrices agree. Affineness gives
the same matrix for every $A_t$. Lemma~\ref{lem:gram} gives equality of
errors. The two endpoint distributions have disjoint nonempty supports,
and all the states are normalized.
\end{proof}

For $n=3$, $k=1$, $\theta=\pi/4$, the probes are explicitly
\begin{align*}
 \ket{D_1}&=(\ket{100}+\ket{010}+\ket{001})/\sqrt3,\\
 \ket{D_2}&=(\ket{110}+\ket{101}+\ket{011})/\sqrt3,\\
 \ket{A_{1/2}}&=(\ket{100}+\ket{010}+\ket{001}
                  +\ket{110}+\ket{101}+\ket{011})/\sqrt6.
\end{align*}
All three final-state Gram matrices equal
$\frac23 I+\frac13J$, with eigenvalues $5/3,2/3,2/3$.
The next lemma computes their error exactly.

\begin{lemma}\label{lem:symmetric}
For $n$ equiprobable normalized pure states with Gram matrix
$G=(1-c)I+cJ\succeq0$, their optimal success equals
\begin{equation}\label{eq:f}
 F_n(c)=\frac{\bigl(\sqrt{1+(n-1)c}+(n-1)\sqrt{1-c}\bigr)^2}{n^2}.
\end{equation}
On $0\leq c\leq1$, this function is nonincreasing.
\end{lemma}
\begin{proof}
Let $B=\sqrt G$ and
$a=B_{ii}=(\sqrt{1+(n-1)c}+(n-1)\sqrt{1-c})/n$,
independent of $i$. Work with the equivalent ensemble $B e_i$ in
$\mathbb C^n$, whose Gram matrix is $G$. The projective measurement
$e_i e_i^*$ attains success $a^2$. For an upper bound take
$Y=aB/n$. Cauchy--Schwarz in the positive semidefinite form defined by
$B$ gives
\[
 |e_i^*Bv|^2\leq(e_i^*Be_i)(v^*Bv)=a\,v^*Bv.
\]
Hence $Y\succeq (Be_i)(Be_i)^*/n$ for every $i$. Any POVM has success
at most $\sum_i\tr(YM_i)=\tr Y=a^2$. This proves the formula even
when $G$ is singular. For $0<c<1$, the derivative of the numerator's
positive square root is
\[
 \frac{n-1}{2}\left(\frac1{\sqrt{1+(n-1)c}}
                         -\frac1{\sqrt{1-c}}\right)\leq0.
\]
Continuity includes both endpoints. This is the square-root measurement
calculation; see~\cite{EF} for its general setting and~\cite[Lemma 2]{ZGH}
for the equal-overlap formula.
\end{proof}

Substitution of $n=3$, $c=1/3$ into~\eqref{eq:f} yields
\begin{equation}\label{eq:exact}
 P(D_1)=P(D_2)=P(A_t)=1-\frac{13+4\sqrt{10}}{27}
 =\frac{14-4\sqrt{10}}{27}\qquad(0\leq t\leq1).
\end{equation}
This refutation has positive error, rather than depending on a zero-error
boundary case.

\section{Weak averaging and the companion conjecture}
Write $P(p)$ for~\eqref{eq:error} with amplitudes $\sqrt{p_x}$.

\begin{theorem}[Weak averaging]\label{thm:convex}
For coefficient-square distributions $p,q$ and $0\leq t\leq1$,
\begin{equation}\label{eq:convex}
 P(tp+(1-t)q)\leq tP(p)+(1-t)P(q).
\end{equation}
Thus equal-error inputs have an average with no larger error; strict
improvement is not guaranteed.
\end{theorem}
\begin{proof}
Let $\phi_i(p)$ and $\phi_i(q)$ be their final states. In the direct sum
of two copies of the Hilbert space put
\[
 \eta_i=\sqrt t\,\phi_i(p)\oplus\sqrt{1-t}\,\phi_i(q).
\]
Their Gram matrix is $tG(p)+(1-t)G(q)=G(tp+(1-t)q)$.
If $M_i$ and $N_i$ are optimal POVMs for the two ensembles, the block
POVM $M_i\oplus N_i$ has success $t(1-P(p))+(1-t)(1-P(q))$
on $\eta_i$. Lemma~\ref{lem:gram} transports this measurement to
the averaged probe, proving~\eqref{eq:convex}. The endpoints are immediate.
The same proof applies to any fixed priors.
\end{proof}

\begin{corollary}[Symmetric optimum]\label{cor:sym}
For equal priors there is an optimal coefficient-square distribution
constant on each Hamming-weight layer.
\end{corollary}
\begin{proof}
A sensor permutation merely relabels outcomes and changes all states by a
common unitary, so it preserves $P$. Average $p$ over all $n!$ sensor
permutations. Repeated use of~\eqref{eq:convex} shows that this average
has error at most $P(p)$, and the permutation orbits are exactly the
weight layers. An optimum exists because the probe sphere and POVM set
are compact and the joint success is continuous. Apply the averaging
argument to such an optimum. This proves existence, not that every
optimum has layer symmetry.
\end{proof}

\begin{corollary}[Conjecture 1]\label{cor:optimal}
Put $m=\lfloor n/2\rfloor$ and
\[
 c_{\min}=1-\frac{4m(n-m)}{n(n-1)}\sin^2\theta,
 \qquad c_*=\max\{0,c_{\min}\}.
\]
The global minimum error over normalized pure probes is $1-F_n(c_*)$.
If $c_{\min}\geq0$, a uniform probe on any layer minimizing $c_k$ is
optimal. This proves Conjecture 1 of~\cite[Section 5, p.~10:10]{ZGH}.
\end{corollary}
\begin{proof}
By Corollary~\ref{cor:sym}, it suffices to consider
$p=\sum_{k=0}^n w_k p^{(k)}$, $w_k\geq0$, $\sum_k w_k=1$.
Its Gram matrix has common overlap $c=\sum_k w_k c_k$.
Since $k(n-k)$ is maximized at the central layers, the smallest $c_k$
is $c_{\min}$, while $c_0=1$.
If $c_{\min}\geq0$, Lemma~\ref{lem:symmetric} maximizes success at
$c=c_{\min}$. If $c_{\min}<0$, mix a central layer with layer zero,
using central-layer weight $1/(1-c_{\min})$. This gives $c=0$ and
success $1$, the universal upper bound. In the source's notation,
$c_k=(R_k+\cos(2\theta)L_k)/(R_k+L_k)$.
Moreover $c_{\min}\geq0$ is exactly
$\theta\in(0,T]\cup[\pi-T,\pi)$, with
$T=\tfrac12\arccos(- (r-1)/r)$ and $r=\lceil n/2\rceil$.
Indeed $\sin^2T=n(n-1)/(4m(n-m))$ for either parity of $n$.
These are precisely the range and minimizers in Conjecture 1.
\end{proof}

The symmetric-probe formula and the construction at $c_*=0$ are
already in~\cite[Section III.B, Eqs.~(16)--(17)]{SLHCM24}, with the
Gram formula and construction present in the 2018 version~\cite{SHCM18}.
Corollary~\ref{cor:sym} supplies the reduction to that class and resolves
the qubit-unitary pure-symmetry question left in~\cite[Section V]{SLHCM24}.

For odd $n$, the two central layers are distinct. Whenever
$c_{\min}\geq0$, every distribution on the segment between them is
globally optimal. In particular the three probes in~\eqref{eq:exact}
are optimal even though their coefficient-squares differ.

\section{Unambiguous discrimination and further consequences}
For unambiguous discrimination a POVM has outcomes $1,\ldots,n,?$;
it never reports $i$ on state $j\ne i$, and its failure probability
is the probability of $?$.

\begin{corollary}\label{cor:ud}
Unambiguous failure is convex in coefficient-squares, and the flat
family of Theorem~\ref{thm:flat} has constant failure for any fixed priors.
For equal priors its global optimum is $c_*=\max\{0,c_{\min}\}$.
A probe attaining the minimum-error optimum in
Corollary~\ref{cor:optimal} also attains this unambiguous optimum.
\end{corollary}
\begin{proof}
Lemma~\ref{lem:gram} preserves every outcome probability, including
zero wrong-answer probabilities. The block measurement in
Theorem~\ref{thm:convex} remains unambiguous when its constituents do;
its inconclusive effect is the direct sum of theirs. This proves
convexity and permits the same permutation averaging.

For a symmetric ensemble with Gram matrix $G=(1-c)I+cJ$, conclusive
probabilities $q_i$ are feasible exactly when
$G-\operatorname{diag}(q_i)\succeq0$ and $q_i\geq0$.
Necessity follows by applying the inconclusive effect to the ensemble:
its Gram contribution is the residual matrix. Conversely, choose vectors
with that residual Gram matrix and map the input vectors to
$\sqrt{q_i}e_i\oplus\beta_i$. Their Gram matrix is $G$, so the
isometry in Lemma~\ref{lem:gram} implements the desired measurement.
Average a feasible diagonal over all outcome permutations. Its mean
$\bar q I$ is feasible, so $\bar q\leq\lambda_{\min}(G)$, and equality
is feasible. Thus the symmetric failure is
\[
 1-\lambda_{\min}(G)=\max\{c,-(n-1)c\}.
\]
For $c\geq0$ this recovers~\cite[Lemma 3]{ZGH}. Every symmetric probe
has $c\in[c_{\min},1]$. If $c_{\min}\geq0$, minimize at $c_{\min}$;
otherwise use the $c=0$ mixture from Corollary~\ref{cor:optimal}.
These are exactly the same probes as for minimum error.
\end{proof}

This proves the unambiguous version of Conjecture 1 proposed in
Section 8.1, item (2), and answers the common-optimal-probe question
in~\cite[Section 9, final sentence, p.~10:23]{ZGH}.
Item (3)'s strict-averaging assertion is refuted even at positive failure:
for $n=3$, $\theta=\pi/4$, every member of the flat family has
failure $1/3$.

\begin{corollary}\label{cor:patterns}
Let $V_a\ket x=v_a(x)\ket x$, $|v_a(x)|=1$, be any finite family
of diagonal unitaries with fixed priors. Minimum error and unambiguous
failure are convex in the coefficient-square distribution. Auxiliary
systems do not improve either optimum.
\end{corollary}
\begin{proof}
The Gram matrix has entries
$G(p)_{ab}=\sum_xp_x\overline{v_a(x)}v_b(x)$, affine in $p$.
The direct-sum argument proves both convexities. For an auxiliary-system
pure probe write $\sum_x\ket x\ket{a_x}$ and put
$p_x=\langle a_x\mid a_x\rangle$. Its output Gram matrix equals
$G(p)$, so the probe $\sum_x\sqrt{p_x}\ket x$ has identical attainable
probabilities. A mixed probe may be purified; discarding the additional
system cannot improve its measurement optimum, giving the same bound.
\end{proof}

The ancillary-system observation for one affected unitary is already
proved in~\cite[Appendix A]{SLHCM24}; the argument above applies
to every diagonal impact pattern and either discrimination objective.
It covers the algebraic reduction for multiple affected sensors or
commuting impact operators mentioned in~\cite[Section 9]{ZGH}, but
does not determine their general optimum. Later multi-event results
are given by Ali--Hillery~\cite{AH}, and perfect trajectory discrimination
by Chin--Chuang~\cite[Theorem 8(b)]{CC}; those are different scopes
from general minimum-error optimization. For unequal priors, Section 8.2's
orthogonality region remains exact when all priors are positive:
zero error is equivalent to orthogonal pure outputs, independently of
the positive prior weights. We do not impose equal priors on the flat-family
or weak-averaging statements.

\begin{corollary}\label{cor:necessity}
The existence of a layer-symmetric optimum does not require every optimum
to be layer-symmetric. For $n=3$, $\theta=\pi/2$, and $0<a<1$,
$p_x=(1+a(2x_1-1))/8$ is a non-layer-symmetric optimum for both criteria.
\end{corollary}
\begin{proof}
This distribution is positive and sums to one. For $i\ne j$ the Gram
entry is $\mathbb E_p[s_i s_j]=0$, because uniform signs have vanishing
nonconstant products, including $\mathbb E[s_1s_i s_j]$.
Hence $G=I$ and both errors are zero. Words of the same weight with
different first bits have different probabilities.
\end{proof}
This universal necessity statement is not posed as a separate conjecture;
Corollary~\ref{cor:necessity} distinguishes it from the existence
question before Conjecture 2.

\begin{corollary}\label{cor:projective}
For the original equal-prior model, restricting the measurement to an
arbitrary global projective measurement leaves both optimal values
unchanged. Unambiguous measurements include an inconclusive outcome.
\end{corollary}
\begin{proof}
Use the common optimal probe with $G=(1-c_*)I+c_*J$.
For minimum error, Lemma~\ref{lem:symmetric} attains $F_n(c_*)$
by the coordinate projectors on the ensemble $\sqrt G e_i$ in
$\mathbb C^n$. The isometry between its span and the physical output
span extends to an isometric embedding of $\mathbb C^n$ in the
$2^n$-dimensional sensor space. Transport the coordinate projectors
and add the remaining orthogonal-complement projector to one outcome.
The resulting effects are mutually orthogonal projectors with the same
success probability.

For unambiguous discrimination use vectors
$\sqrt{1-c_*}e_i+\sqrt{c_*}e_?$ in $\mathbb C^{n+1}$, which
also have Gram matrix $G$. Since $2^n\geq n+1$ for $n\geq3$,
the same extension embeds the coordinate projectors in the physical
space. Assign the remaining complement to the inconclusive outcome.
This projective measurement never gives a wrong conclusive answer
and has failure $c_*$.
\end{proof}
Thus the global-projective part of~\cite[Section 9]{ZGH}'s
measurement-restriction question has the same optimum. The
minimum-error square-root measurement for symmetric probes is already
in~\cite{EF,SLHCM24}. Local-measurement restrictions and noisy-channel
optimization (Section 8.3) remain outside the present conclusions.

\section{Coverage and mathematical evidence}

\begin{center}\small
\begin{tabular}{@{}p{.40\linewidth}p{.53\linewidth}@{}}
\toprule
Source item or mathematical statement&Coverage\\\midrule
Complementary-layer Gram formula&Reproved; prior proof Skotiniotis
et al.\ (2018, Eq.~(8); 2024, Eq.~(14)).\\
Section 5, Conjecture 2, p.~10:11&Refuted-and-corrected:
Theorems~\ref{thm:flat} and~\ref{thm:convex}.\\
Section 5, symmetry-of-coefficients question before Conjecture 2&Answered: yes
(existence):
Corollary~\ref{cor:sym}.\\
Skotiniotis et al.\ (2024), Section V, qubit-unitary pure-symmetry question&Answered:
yes, Corollary~\ref{cor:sym}.\\
Section 5, Conjecture 1, p.~10:10&Proved without symmetry assumption:
Corollary~\ref{cor:optimal}; symmetric calculation reproved, prior
Skotiniotis et al.\ (2018; 2024).\\
Section 4, Theorem 2 and Corollary 1&Reproved; prior proof
Zhan--Gupta--Hillery (2023 preprint; 2024), with constructive
sufficiency in Skotiniotis et al.\ (2018); later independent
Chin--Chuang (2024 preprint; 2025).\\
Section 8.1, items (2)--(3)&Conjecture 1 variant proved;
strict averaging refuted-and-corrected: Corollary~\ref{cor:ud}.\\
Section 8.2&Complementary-layer equality and weak averaging hold for
fixed priors; general unequal-prior optimum not covered.\\
Section 8.3&Noisy-channel optimization not covered.\\
Section 9, final question, p.~10:23&Answered: yes, common optimal probe,
Corollary~\ref{cor:ud}.\\
Section 9, multiple impacts&General optimum not covered;
diagonal-pattern reduction proved, Corollary~\ref{cor:patterns}.\\
Section 9, restricted measurements&Answered for global projective measurements:
both optima are preserved, Corollary~\ref{cor:projective}; local restrictions not covered.\\
Universal symmetry of every optimum&Not posed as a separate conjecture;
fails for Corollary~\ref{cor:necessity}'s example.\\\bottomrule
\end{tabular}
\end{center}
The arguments above are written proofs. The accompanying \texttt{verify.py}
checks exact layer counts for $3\leq n\leq8$, the three-sensor Gram matrices,
and a matching measurement/upper-bound certificate over
$\mathbb Q(\sqrt{10})$ for~\eqref{eq:exact}.
The additional rational checks cover the nonsymmetric three-sensor
family at thirteen choices of $a$, equal-Gram ancillary probes and the
residual-Gram characterization for unambiguous discrimination.
In the strict-averaging Lean file linked below, the declaration
\texttt{result} proves the negation
of the universal strict claim using $n=3$, $\theta=\pi/4$ and the two
optimal layers: its comparison is that the average cannot have lower
error than an optimal endpoint. Equality along the whole segment and
the radical value in~\eqref{eq:exact} are written proofs here, not
conclusions of that declaration. In the companion Dicke-optimality file,
\texttt{result} compares each proposed
minimizing-layer probe against every normalized pure competitor and every
global equal-prior POVM in the angular range of Conjecture 1.
The corresponding fixed-revision sources are
\href{https://github.com/the-omega-institute/trureturing/blob/15070fa5f71c4237050acc7ebe058958cbc04e6e/D5/S3/Quantum/Measurements/DetectorSensorStrictAveragingRefutation.lean}{the strict-averaging declaration}
and
\href{https://github.com/the-omega-institute/trureturing/blob/15070fa5f71c4237050acc7ebe058958cbc04e6e/D5/S3/Quantum/Measurement/DetectorSensorDickeOptimality.lean}{the Dicke-optimality declaration}.
The general flat-family identity, weak-averaging
proof, unambiguous and diagonal-pattern extensions are established above
by written proofs; the rational and quadratic-field checks have the
finite scope specified here.

\begin{thebibliography}{9}
\bibitem{ZGH} C.~Zhan, H.~Gupta and M.~Hillery,
\emph{Optimizing Initial State of Detector Sensors in Quantum Sensor Networks},
ACM Transactions on Quantum Computing \textbf{5}(2), Article 10 (2024), 1--25.
\href{https://doi.org/10.1145/3655028}{doi:10.1145/3655028}.
\href{https://arxiv.org/abs/2306.17401v6}{arXiv:2306.17401v6}.
\bibitem{EF} Y.~C.~Eldar and G.~D.~Forney, Jr.,
\emph{On quantum detection and the square-root measurement},
IEEE Transactions on Information Theory \textbf{47}(3) (2001), 858--872.
\href{https://doi.org/10.1109/18.915636}{doi:10.1109/18.915636}.
\href{https://arxiv.org/abs/quant-ph/0005132v2}{arXiv:quant-ph/0005132v2}.
\bibitem{HGZ} M.~Hillery, H.~Gupta and C.~Zhan,
\emph{Discrete outcome quantum sensor networks},
Physical Review A \textbf{107} (2023), 012435.
\href{https://doi.org/10.1103/PhysRevA.107.012435}{doi:10.1103/PhysRevA.107.012435}.
\href{https://arxiv.org/abs/2210.17254}{arXiv:2210.17254}.
\bibitem{SHCM18} M.~Skotiniotis, R.~Hotz, J.~Calsamiglia and
R.~Mu\~noz-Tapia, \emph{Identification of malfunctioning quantum devices},
arXiv:1808.02729v1 (August 8, 2018), Eqs.~(8)--(11).
\href{https://arxiv.org/abs/1808.02729v1}{2018 version}.
\bibitem{SLHCM24} M.~Skotiniotis, S.~Llorens, R.~Hotz, J.~Calsamiglia
and R.~Mu\~noz-Tapia, \emph{Identification of malfunctioning quantum devices},
Physical Review Research \textbf{6} (2024), 033329.
\href{https://doi.org/10.1103/PhysRevResearch.6.033329}{doi:10.1103/PhysRevResearch.6.033329}.
\href{https://arxiv.org/abs/1808.02729v3}{arXiv:1808.02729v3},
Section III.B and Appendices A--B.
\bibitem{AH} N.~Ali and M.~Hillery, \emph{Discrete-outcome sensor
networks. II. Multiple detection events and grouping detectors},
Physical Review A \textbf{110} (2024), 012619.
\href{https://doi.org/10.1103/PhysRevA.110.012619}{doi:10.1103/PhysRevA.110.012619}.
\href{https://arxiv.org/abs/2407.20435}{arXiv:2407.20435}.
\bibitem{CC} Z.~E.~Chin and I.~L.~Chuang,
\emph{Quantum trajectory sensing problem and its solution},
Physical Review A \textbf{111} (2025), 052447.
\href{https://doi.org/10.1103/PhysRevA.111.052447}{doi:10.1103/PhysRevA.111.052447}.
\href{https://arxiv.org/abs/2410.00893v1}{arXiv:2410.00893v1}
(October 1, 2024), Theorem 8(b).
\end{thebibliography}
\end{document}
